% !PRINTED 258
\weberpage{258}

In the case $n=1$, $p_{,}=5$, however, we have in $G$ by (2) % [sic] odd subscript comma printed after p
and 3) the substitution

\[
\begin{pmatrix} 1, & 1 \\ 2, & -2 \end{pmatrix}
\begin{pmatrix} 0, & 1 \\ -1, & 0 \end{pmatrix}
\begin{pmatrix} -2, & -1 \\ -2, & 1 \end{pmatrix}
\begin{pmatrix} 0, & 1 \\ -1, & 0 \end{pmatrix}
=
\begin{pmatrix} -2, & 0 \\ 0, & 2 \end{pmatrix},
\]
and consequently also
\weq{12}{
\begin{pmatrix} 1, & 0 \\ -\xi, & 1 \end{pmatrix}
\begin{pmatrix} -2, & 0 \\ 0, & 2 \end{pmatrix}
\begin{pmatrix} 1, & 0 \\ \xi, & 1 \end{pmatrix}
\begin{pmatrix} 2, & 0 \\ 0, & -2 \end{pmatrix}
=
\begin{pmatrix} 1, & 0 \\ -2\xi, & 1 \end{pmatrix},
}
hence, since $\xi$ and consequently also $-2\xi$ is arbitrary, again the substitutions (10) and with them therefore all the generators of the group $E$,
\gesperrt{and so in all these cases $G$ is identical with $E$ and
consequently the group $E$ is simple}.

The case $n=1$, $p=3$, however, presents the second genuine
exception. In this case the group $E$ is of degree $12$
and it has the normal divisor of degree $4$

\[
G =
\begin{pmatrix} 1, & 0 \\ 0, & 1 \end{pmatrix},
\begin{pmatrix} 0, & 1 \\ -1, & 0 \end{pmatrix},
\begin{pmatrix} -1, & 1 \\ 1, & 1 \end{pmatrix},
\begin{pmatrix} 1, & 1 \\ 1, & -1 \end{pmatrix}.
\]

To convince oneself of this, one need only compare the two
subsidiary groups:
\[
\begin{aligned}
\begin{pmatrix} 1, & 0 \\ 1, & 1 \end{pmatrix} G &=
\begin{pmatrix} 1, & 0 \\ 1, & 1 \end{pmatrix},
\begin{pmatrix} 0, & 1 \\ -1, & 1 \end{pmatrix},
\begin{pmatrix} -1, & 1 \\ 0, & -1 \end{pmatrix},
\begin{pmatrix} 1, & 1 \\ -1, & 0 \end{pmatrix} \\
\begin{pmatrix} 1, & 0 \\ -1, & 1 \end{pmatrix} G &=
\begin{pmatrix} 1, & 0 \\ -1, & 1 \end{pmatrix},
\begin{pmatrix} 0, & 1 \\ -1, & -1 \end{pmatrix},
\begin{pmatrix} -1, & 1 \\ -1, & 0 \end{pmatrix},
\begin{pmatrix} 1, & 1 \\ 0, & 1 \end{pmatrix}
\end{aligned}
\]
with $G\begin{pmatrix} 1, & 0 \\ 1, & 1 \end{pmatrix}$, $G\begin{pmatrix} 1, & 0 \\ -1, & 1 \end{pmatrix}$.

With this it is therefore proved in general that, apart from the two
cases $n=1$, $p=2$ and $n=1$, $p=3$, the group $E$
is simple. For the first cases one obtains for the degrees $g$
of these simple groups:

\[
\begin{array}{lll@{\qquad}lll}
n=1, & p=\ 5, & g=\ \ 60 & n=2, & p=2, & g=\ \ 60 \\
     & p=\ 7, & g=\ 168  &      & p=3, & g=\ 360 \\
     & p=11,  & g=\ 660  &      & p=5, & g=7800 \\
     & p=13,  & g=1092   & n=3, & p=2, & g=\ 504 \\
     &        &          &      & p=3, & g=9828 \\
     &        &          & n=4, & p=2, & g=4080^{1)}.
\end{array}
\]

\rule{4cm}{0.4pt}

$^{1)}$ See Bulletin of the New York math. Society. October 1893. In
the report on the mathematical congress in Chicago there are
two notes on these groups by \gesperrt{Cole} and \gesperrt{Moore}.
